\appendix
%  \appendix switches \chaptertitlename to "Appendix" and the numbering to
%  letters, so these render as "Appendix A / Title" in the same style as the
%  chapters.

% \chapter{Additional Thesis Details}
% \clearpage

\chapter{Additional Survey Details}

\section{Comparison with Prior Surveys}
\label{app:survey-comparison}
\begin{table*}[h]
\centering
\setstretch{1}
\footnotesize
\setlength{\tabcolsep}{5pt}
\renewcommand{\arraystretch}{1.15}
\newcolumntype{L}[1]{>{\raggedright\arraybackslash\hsize=#1\hsize}X}
\begin{tabularx}{\textwidth}{@{}L{1.75} L{0.5} L{1.0} L{0.75} L{1.0}@{}}
\toprule
\textbf{Survey} & \textbf{Venue} & \textbf{Scope} & \textbf{Multi-Page Treatment} & \textbf{Organising Framing} \\
\midrule
A Survey of Deep Learning Approaches for OCR and Document Understanding~\citeauthor{subramani2020survey} & arXiv 2020 & Pre-MLLM OCR pipelines & Not covered & Pipeline stage (detection, recognition, layout) \\
Deep Learning Approaches for Information Extraction from Visually Rich Documents: Datasets, Challenges and Methods~\citeauthor{gbada2025deep} & IJDAR 2025 & Information extraction over VRDs & Not covered & Extraction task and method \\
Deep Learning based Visually Rich Document Content Understanding: A Survey~\citeauthor{ding2026deep} & AI Review 2026 & Content understanding over VRDs & Not covered & Task formulation, architecture, pretraining \\
A Survey on MLLM-based Visually Rich Document Understanding: Methods, Challenges, and Emerging Trends~\citeauthor{ding2026survey} & ACL 2026 & MLLM-era VRDU & In passing & OCR dependency; architecture, training, benchmarks \\
Scaling Beyond Context: A Survey of Multimodal Retrieval-Augmented Generation for Document Understanding~\citeauthor{gao2026scaling} & ACL 2026 & Multimodal RAG for long documents & Partial (RAG techniques only) & Retrieval-augmented generation \\
\midrule
\textbf{This survey} & EMNLP 2026 & \textbf{Multi-page VRDU} & \textbf{Distinct problem} & \textbf{Evidence management} \\
\bottomrule
\end{tabularx}
\caption{Comparison of this survey with prior VRDU surveys. Prior surveys either predate the multi-page literature or address it only within a broader single-page or extraction-focused scope. This survey is the first to treat multi-page VRDU as a distinct problem organised around evidence management.}
\label{tab:survey-comparison}
\end{table*}

Table~\ref{tab:survey-comparison} contrasts this survey with the five prior VRDU surveys closest to our scope. The comparison is along three axes. \textbf{Scope} records the era and problem each survey targets, from pre-MLLM OCR pipelines to the recent MLLM and retrieval-augmented literature. \textbf{Multi-page treatment} records whether multi-page understanding is treated as a distinct problem with its own design axes, addressed only in passing within a broader scope, or left out. \textbf{Organising framing} records the lens each survey uses to structure the field. Prior surveys either predate the multi-page literature or treat it as one topic within a single-page or extraction-focused scope. This survey is the first to define MP-VRDU as a distinct problem, to organise it around evidence management, and to cover the four architectural families, modality pressures, inference-time and training strategies, and datasets specific to the multi-page setting.

\section{Full Survey Methodology}
\label{app:sv_methodology}

\subsection{Search keywords.}
Queries combined a multi-page cluster with a task or method cluster. Multi-page terms included ``multi-page document'', ``long document'', ``long-context document'', ``cross-page reasoning'', ``page-level retrieval'', ``document-scale evidence'', and ``hundred-page''. Task and method terms included ``document visual question answering'', ``visually rich document understanding'', ``document VQA'', ``document agent'', ``visual document retrieval'', ``document RAG'', ``hierarchical document transformer'', and ``long-context vision-language model''. Benchmark-name probes (``MP-DocVQA'', ``MMLongBench-Doc'', ``LongDocURL'', ``DUDE'', ``SlideVQA'', ``M3DocVQA'', ``DocBench'', ``M-LongDoc'') were used as anchors for forward-citation sweeps to capture systems that report against multi-page evaluation suites.

\subsection{Inclusion and exclusion criteria.}
A paper was retained if it (i) proposed a multi-page model, framework, system, or pipeline rather than a benchmark or dataset alone, and (ii) reported quantitative results on at least one multi-page benchmark, with cross-page evidence handling as an explicit design contribution. Excluded categories were single-page-only systems benchmarked solely on DocVQA, InfographicVQA, ChartQA, FUNSD, or CORD; survey, position, and opinion papers; generic OCR or table-extraction tools without VQA or reasoning evaluation; long-context LLM benchmarks where documents are one input type among many; and dataset-only contributions, which are catalogued separately in \S\ref{sec:datasets}. Peer-reviewed venue publications and arXiv preprints were assessed against the same substantive criteria, since recent MP-VRDU progress has appeared predominantly in preprint form; venue status is recorded per entry in Table~\ref{tab:model_survey_arch}. Table~\ref{tab:survey_corpus} lists the surveyed systems grouped by architectural family, with the multi-page contribution on which each was included.

\subsection{Knowledge extraction and verification.}
Each retained paper was converted to markdown and processed through LLM-assisted key-information extraction with evidence-span location, populating per-model fields covering architectural family, OCR dependency, vision encoder, LLM backbone, adaptors and projectors, page-rendering resolution, training stages, datasets used per stage, and reported scores on each benchmark. Every extracted field was manually verified against the source paper before being added.

\subsection{Dataset catalogue.}
The dataset catalogue in \S\ref{sec:datasets} was built primarily from the pretraining, fine-tuning, and benchmarking datasets named by every paper in the surveyed corpus, supplemented with a small set of datasets the corpus does not yet adopt but that cover MP-VRDU-relevant settings such as long scientific documents, multi-page tables, and cross-page reasoning suites, included to surface datasets whose adoption may strengthen future cross-system comparison.

\begin{table*}[p]
\centering
% The document runs at \setstretch{1.1}. A 50-row table does not need that extra
% leading, and giving it back is most of what makes this fit on one page.
\setstretch{1}
\scriptsize
\captionsetup{font=small}
% A wider description column is what actually buys the page: at 0.74 most rows
% wrap to two lines, at 0.78 most do not, which is worth ~140pt of height. The
% table then overruns \textwidth by 6%, and the adjustbox below scales that back.
\setlength{\tabcolsep}{3pt}
\renewcommand{\arraystretch}{1.0}
% Defined here as well as in survey/models_summary.tex, so this table does not
% depend on that one having been input first.
\definecolor{sectionbg}{gray}{0.88}
% Backstop: only scales if the rows still overrun, and leaves the caption alone.
\begin{adjustbox}{max totalsize={\textwidth}{0.93\textheight},center}
\begin{tabular}{@{}l l >{\raggedright\arraybackslash}p{0.78\textwidth}@{}}
\toprule
\textbf{Model} & \textbf{Venue} & \textbf{Multi-Page Contribution} \\
\midrule

\rowcolor{sectionbg}
\multicolumn{3}{l}{\textbf{Page-to-Document Encoding Models}} \\
\midrule
  {Arctic-TILT}~(\citeyear{borchmann2025arctic}) & ACL & Layer-wise tensor-product text-image fusion with chunked sparse attention, serving up to 389k tokens on a single GPU \\
  {GRAM}~(\citeyear{blau2024gram}) & CVPR & Interleaves page-level and document-level attention sub-layers with learnable global Doc tokens, avoiding quadratic concatenation cost \\
  {Hi-VT5}~(\citeyear{tito2023hierarchical}) & PR & Learnable [PAGE] tokens summarise each page; decoder consumes the concatenated page tokens, with answer-page prediction for explainability \\
  {RM-T5}~(\citeyear{dong2024multi}) & DAS & Recurrent memory tokens propagate context page-by-page; per-page memories concatenated for decoding \\

\midrule[0.4pt]

\rowcolor{sectionbg}
\multicolumn{3}{l}{\textbf{MLLM-Centric Adaptation Models}} \\
\midrule
  {CoR}~(\citeyear{guo2026end}) & preprint & Four-stage Chain-of-Reading paradigm with a caption-reconstruction objective for figure and table grounding \\
  {Docopilot}~(\citeyear{duan2025docopilot}) & CVPR & Native multi-page MLLM trained on Doc-750K cross-page proxy tasks with context-extension engineering \\
  {DocR1}~(\citeyear{xiong2026docr1}) & AAAI & EviGRPO adds an evidence-page F1 reward over structured think/evidence/answer outputs \\
  {DocSeeker}~(\citeyear{yan2026docseeker}) & CVPR & Analysis-Localization-Reasoning paradigm with page-identifier grounding and evidence-guided resolution allocation \\
  {DocSLM}~(\citeyear{hannan2026docslm}) & CVPR & Hierarchical compressor fixes 576 tokens per page; streaming segments with entropy-based abstention for edge deployment \\
  {InstructDoc}~(\citeyear{tanaka2024instructdoc}) & AAAI & Document-former fuses OCR, layout, and visual features into learnable tokens, with pages encoded in parallel \\
  {LayTokenLLM}~(\citeyear{zhu2025simple}) & CVPR & Per-segment layout token shares its segment's position ID, adding layout without consuming context \\
  {Leopard}~(\citeyear{jia2025leopard}) & TMLR & Adaptive high-resolution multi-image encoding under a global token budget with pixel-shuffle compression \\
  {mPLUG-DocOwl2}~(\citeyear{hu2025mplug}) & ACL & DocCompressor cross-attends global low-resolution queries to sub-image features, compressing each page to 324 tokens \\
  {Texthawk2}~(\citeyear{yu2024texthawk2}) & preprint & Two-step resampling and rearrangement for 16$\times$ visual-token compression with a coordinate-token detection head \\
  {URaG}~(\citeyear{shi2026urag}) & AAAI & Cross-modal retrieval module at an early Transformer layer discards non-retained pages before deeper layers \\

\midrule[0.4pt]

\rowcolor{sectionbg}
\multicolumn{3}{l}{\textbf{Retrieval-Augmented Pipelines}} \\
\midrule
  {AVIR}~(\citeyear{li2026avir}) & MMAsia & Adaptive page selector uses score-distribution clustering to set per-query retrieval depth \\
  {CREAM}~(\citeyear{zhang2024cream}) & ACM & Coarse-to-fine retrieval with iterative reranking, paired with attention-pooled multi-page visual features \\
  {DFVC}~(\citeyear{qian2026accurate}) & MDPI & Adapter predicts gating weights over neighbouring-page embeddings to fuse cross-page context into page representations \\
  {DREAM}~(\citeyear{zhang2025dream}) & MM & Learning-to-rank fusion of embedding similarity and keyword-guided MLLM scores, with cross-page mixture-of-experts routing \\
  {KGP}~(\citeyear{wang2024knowledge}) & AAAI & Knowledge graph over passage and structural nodes traversed by an LLM agent to gather supporting facts \\
  {M2RAG}~(\citeyear{duan2025m2rag}) & NeurIPS & Dual-tower sparse and dense retrieval with text, visual, and fusion agents performing consistency checking \\
  {M3DocRAG}~(\citeyear{cho2024m3docrag}) & preprint & ColPali late-interaction page retrieval with approximate indexing, scaling to 40k+ pages \\
  {MDocAgent}~(\citeyear{han2025mdocagent}) & preprint & Five-stage pipeline with parallel text and image retrieval and specialised modality agents \\
  {MHier-RAG}~(\citeyear{gong2026mhier}) & preprint & Combines flattened in-page indexing with a recursively summarised cross-page tree for multi-granularity retrieval \\
  {MLDocRAG}~(\citeyear{zhang2026mldocrag}) & preprint & Chunk-Query Graph of LVLM-generated queries links chunks for retrieval by query-to-query similarity \\
  {MM-R5}~(\citeyear{xu2025mm}) & AAAI & Two-stage reranker training with structured reasoning chains and rank-discounted GRPO rewards \\
  {MoLoRAG}~(\citeyear{wu2025molorag}) & EMNLP & Traverses a page graph combining semantic and VLM-judged logical relevance to reach weakly-similar evidence \\
  {MultiDocFusion}~(\citeyear{shin2026multidocfusion}) & ACL & LLM-constructed document hierarchy tree drives depth-first hierarchical chunking \\
  {PDF-WuKong}~(\citeyear{xie2026pdf}) & IJCV & End-to-end sparse sampler jointly trained with the LLM selects top-K text and image evidence \\
  {RAG-DocVQA}~(\citeyear{pez2025enhancing}) & ICDAR & Systematic comparison of textual and visual RAG variants for multi-page DocVQA \\
  {Self-Attention Scoring}~(\citeyear{kang2024multi}) & ICDAR & Self-attention scoring module ranks independently encoded pages, extending evaluation to 793-page documents \\
  {SLEUTH}~(\citeyear{liu2025resolving}) & CVPR & Training-free clue discovery and page screening keep effective context fixed as retrieval depth grows \\
  {SV-RAG}~(\citeyear{chen2025sv}) & ICLR & Dual LoRA adapters over one MLLM backbone serve both late-interaction page retrieval and answering \\
  {VDocRAG}~(\citeyear{tanaka2025vdocrag}) & CVPR & Compresses each page image to a single dense embedding via self-supervised compression pretraining tasks \\
  {VisDoMRAG}~(\citeyear{suri2025visdom}) & NAACL & Parallel textual and visual RAG chains reconciled by consistency-constrained modality fusion \\

\midrule[0.4pt]

\rowcolor{sectionbg}
\multicolumn{3}{l}{\textbf{Adaptive-Trajectory Pipelines}} \\
\midrule
  {DMAP}~(\citeyear{fu2026dmap}) & WWW & Hierarchical document map traversed by tri-path retrieval with a reflective loop for insufficient evidence \\
  {Doc-V*}~(\citeyear{zheng2026doc}) & ACL & Thumbnail overview then alternating retrieval and high-resolution page fetch, trained by SFT then GRPO \\
  {DocAgent}~(\citeyear{sun2025docagent}) & EMNLP & Hierarchical XML outline navigated by an actor agent with five tools, plus a reviewer and persistent memory bank \\
  {DocDancer}~(\citeyear{zhang2026docdancer}) & preprint & End-to-end trained agent with only search and read tools, learned from synthesised exploration trajectories \\
  {DocLens}~(\citeyear{zhu2026doclens}) & ACL & Page navigator and element localiser zoom to fine-grained evidence; sampling-adjudication selects among candidates \\
  {DocReact}~(\citeyear{wu2025doc}) & ACL & Frames iterative retrieval as information-gain maximisation, halting when residual entropy stops decreasing \\
  {MACT}~(\citeyear{yu2025visual}) & preprint & Four specialised agents with independent judgment routing and agent-wise adaptive test-time scaling \\
  {MM-Doc-R1}~(\citeyear{lin2026mm}) & ACL & Planner-seeker-answer workflow trained by similarity-weighted policy optimisation for multi-turn RL \\
  {SimpleDoc}~(\citeyear{jain2025simpledoc}) & EMNLP & Dual-cue retrieval over embeddings and page summaries, with a memory-carrying reasoner that reissues refined queries \\
  {ViDoRAG}~(\citeyear{wang2025vidorag}) & EMNLP & Gaussian-mixture adaptive top-K retrieval with seeker, inspector, and answer agents over multi-scale views \\
  {VRAG-RL}~(\citeyear{wang2025vrag}) & NeurIPS & Visual perception action space lets the agent crop and zoom regions, trained with retrieval-efficiency rewards \\

\bottomrule
\end{tabular}
\end{adjustbox}
\caption{The 46 surveyed MP-VRDU systems. Grouped by architectural family, with the multi-page contribution motivating each system's inclusion. Venue indicates the publication venue at time of writing.}
\label{tab:survey_corpus}
\end{table*}

\section{Additional Resources}

\subsection{Reported Performance by Architecture}
\label{app:performance}

\input{appendix/survey/table-performance-arch}

Table~\ref{tab:model_performance_arch} summarises reported scores on the five most-adopted multi-page benchmarks, grouped by architectural family. It is included for coverage rather than ranking. The matrix is sparse and metric conventions diverge across MP-DocVQA, DUDE, LongDocURL, MMLongBench-Doc, and SlideVQA, so scores are not comparable across columns; for the same benchmark, differences in backbone scale, proprietary-LLM dependency, evidence budget, and the training-versus-zero-shot distinction mean that a higher score does not by itself indicate a better architecture. Several of the strongest systems orchestrate closed-source frontier backbones of undisclosed scale, so their scores cannot be attributed to architecture alone. Read with these caveats, the table supports one robust cross-family observation: coverage tracks cost structure, with page-to-document and MLLM-centric systems confined to mid-length benchmarks and selective-evidence families dominating the long-document frontier.

\subsection{General Domain MLLM for MP-VRDU}
\label{app:general-MLLM}
Frontier general-domain MLLMs handle short multi-page inputs well when the document fits in context, but the multi-page setting introduces three pressures absent in single-page VRDU: context-window overrun, attention dilution over sparsely distributed evidence~\cite{liu2024lost}, and a per-page reading fidelity versus whole-document coverage trade-off under a fixed token budget. Table~\ref{tab:model_performance_arch} shows that even the strongest frontier MLLMs trail specialised MP-VRDU systems by a wide margin on MMLongBench-Doc and LongDocURL, motivating the specialised architectures examined in \S\ref{sec:architecture}.

\subsection{Model Components Details}
\label{app:components}
Table~\ref{tab:model_components} lists the concrete components of each surveyed system, organised by architectural family. The four families assemble these parts in characteristically different ways. Page-to-document transformers pair smaller LLM backbones (T5 variants, DocFormerv2) with custom layout-aware projectors. MLLM-centric adaptation models reuse standard vision-language stacks such as Qwen2.5-VL, mPLUG-Owl2, and SigLIP variants, and update only light projection or compression modules. Retriever-augmented pipelines decouple a frozen retriever (ColPali, ColQwen, VisRAG) from a frozen LVLM and train only the scorer, fuser, or reranker. Adaptive-trajectory pipelines mostly orchestrate frozen frontier VLMs, with parameter updates restricted to a few systems that fine-tune the controller via SFT or GRPO.

\input{appendix/survey/table-components-detailed}

\clearpage
\chapter{Additional Attribution Details}

\section{Extended Result Tables}
\label{app:extended_findings}

\subsection{Representation Ladder by Document Domain}
\label{app:ladder_domain}
Table~\ref{tab:ladder_domain_resolution} breaks the representation ladder of Table~\ref{tab:ladder_source} down by document domain, and sweeps the two image-carrying rungs across render resolution. Domain rows partition the pool, so each question is counted exactly once. The ordering of rungs is stable within a row, with TLV winning in six of the seven domains, but variation across domains is considerably wider than variation across rungs: at TLV accuracy ranges from 42.9\% on academic papers to 73.2\% on tutorial and workshop material. Resolution matters most where the image is the only channel, with pooled accuracy at V rising 12.4 points from low to high against 5.5 points at TLV, and the largest gains falling on the domains dominated by dense tables and figures.
% GENERATED by tables/scripts/ -- do not edit by hand, re-run the script.
% Numbers read from results/tables/mmlongbench/headline_full_pool.csv (the ladder by
% domain), faithfulness_answerable_accuracy.csv (its All row), resolution.csv (the
% low/high wings by domain) and resolution_source.csv (their All row).
% Ladder by document domain at medium resolution, with TLV and V swept across low and
% high. Teal is accuracy on the SAME fixed [14,74] scale as findings/ladder.tex, capped
% at 60%; bold marks each row's winning rung within the medium ladder, ties included.
% The last column is NOT an accuracy: it is V minus TLV at medium, a point difference,
% so it takes its own warm scale over [0,13] points.
\definecolor{accteal}{HTML}{1BAF7A}
\definecolor{recwarm}{HTML}{D95926}
\newcommand{\ldashade}[1]{%
  \pgfmathtruncatemacro{\shadeval}{max(0,min(60,(#1-14)/(74-14)*65))}%
  \edef\tmpcol{accteal!\shadeval}%
  \expandafter\cellcolor\expandafter{\tmpcol}%
}
\newcommand{\ldacell}[1]{\ldashade{#1}#1}
\newcommand{\ldabest}[1]{\ldashade{#1}\textbf{#1}}
\newcommand{\lddshade}[1]{%
  \pgfmathtruncatemacro{\shadeval}{max(0,min(55,(#1-0)/(13-0)*55))}%
  \edef\tmpcol{recwarm!\shadeval}%
  \expandafter\cellcolor\expandafter{\tmpcol}%
}
\newcommand{\lddcell}[1]{\lddshade{#1}$-$#1}
\newcommand{\ldugain}[1]{\lddshade{#1}$+$#1}

\begin{table}[ht]
\centering
\small
\setlength{\tabcolsep}{5pt}
\renewcommand{\arraystretch}{1.12}
\adjustbox{max width=\textwidth}{%
\begin{tabular}{@{} l rrrr @{\hskip 12pt} rr @{\hskip 6pt} rr @{\hskip 14pt} r @{}}
\toprule
 & \multicolumn{4}{c}{\textbf{Ladder @ med (\%)}} & \multicolumn{2}{c}{\textbf{TLV}} & \multicolumn{2}{c}{\textbf{V}} & \\
\cmidrule(lr){2-5}\cmidrule(lr){6-7}\cmidrule(lr){8-9}
\textbf{Domain} & \textbf{T} & \textbf{TL} & \textbf{TLV} & \textbf{V} & \textbf{low} & \textbf{high} & \textbf{low} & \textbf{high} & \textbf{$\Delta$V} \\
\midrule
Academic paper       & \ldacell{36.4}   & \ldacell{33.1}   & \ldabest{42.9}   & \ldacell{33.8}   & \ldacell{37.0}   & \ldacell{46.8}   & \ldacell{19.5}   & \ldacell{38.3}   & \lddcell{9.1} \\
Admin / Industry     & \ldacell{50.0}   & \ldacell{56.2}   & \ldabest{67.2}   & \ldacell{54.7}   & \ldacell{64.1}   & \ldacell{70.3}   & \ldacell{42.2}   & \ldacell{64.1}   & \lddcell{12.5} \\
Brochure             & \ldacell{28.6}   & \ldacell{32.5}   & \ldabest{45.5}   & \ldacell{40.3}   & \ldacell{40.3}   & \ldacell{48.1}   & \ldacell{39.0}   & \ldacell{49.4}   & \lddcell{5.2} \\
Financial report     & \ldacell{49.1}   & \ldabest{52.8}   & \ldacell{51.9}   & \ldacell{38.9}   & \ldacell{54.6}   & \ldacell{53.7}   & \ldacell{21.3}   & \ldacell{46.3}   & \lddcell{13.0} \\
Guidebook            & \ldacell{35.8}   & \ldacell{40.0}   & \ldabest{51.7}   & \ldacell{45.0}   & \ldacell{48.3}   & \ldacell{51.7}   & \ldacell{38.3}   & \ldacell{47.5}   & \lddcell{6.7} \\
Research / Intro     & \ldacell{22.6}   & \ldacell{27.4}   & \ldabest{47.6}   & \ldacell{45.3}   & \ldacell{42.9}   & \ldacell{50.0}   & \ldacell{38.2}   & \ldacell{47.6}   & \lddcell{2.3} \\
Tutorial / Workshop  & \ldacell{14.3}   & \ldacell{48.2}   & \ldabest{73.2}   & \ldacell{67.9}   & \ldacell{71.4}   & \ldacell{74.1}   & \ldacell{74.1}   & \ldacell{70.5}   & \lddcell{5.3} \\
\midrule
\textbf{All}         & \ldacell{31.9}   & \ldacell{38.8}   & \ldabest{52.5}   & \ldacell{45.6}   & \ldacell{49.2}   & \ldacell{54.7}   & \ldacell{37.8}   & \ldacell{50.2}   & \lddcell{6.9} \\
\bottomrule
\end{tabular}%
}
\caption{\small Representation ladder by document domain and render resolution. Teal deepens with accuracy and bold marks each row's winning rung; the warm strip deepens with how far vision alone falls below the combined rung. Resolution lifts the domains whose evidence is pixel-limited and leaves the text-heavy ones where they were.}
\label{tab:ladder_domain_resolution}
\end{table}

% {\footnotesize\raggedright
% Accuracy (\%) on oracle pages, Qwen3-VL-8B bf16, PaddleOCR-VL, answerable pool
% ($n{=}847$). Teal deepens with accuracy on the same fixed scale as
% Table~\ref{tab:ladder_source}, and bold marks each row's winning rung within the
% medium-resolution ladder, both sides of a tie included. The \textbf{Ladder @ med} block
% is the four rungs at medium resolution: text alone (\textbf{T}), parser layout
% (\textbf{TL}), parser layout with the page image (\textbf{TLV}), and the image alone
% (\textbf{V}). The \textbf{TLV} and \textbf{V} blocks give those two rungs at low and
% high resolution. \textbf{$\Delta$V} is V minus TLV at medium: a point difference and not
% an accuracy, so it carries its own warm shading, deepening with the accuracy the text
% channel recovers over vision alone.
% \textbf{Domain rows partition the pool}, unlike the evidence-source rows of
% Table~\ref{tab:ladder_source}, so each question is counted exactly once and the
% \textbf{All} row is the same pool read whole.
% \textbf{The medium column and the low/high wings come from different runs.} The ladder
% block is the ladder run and the wings are the resolution sweep, both on the same 847
% questions; they agree pooled and differ by a point or two within a domain, so read a
% wing against its own block rather than differencing it against the medium column.\par}

\subsection{Parser Fidelity by Evidence Source}
\label{app:parser_fidelity}
Table~\ref{tab:parser_source_scan} splits each parser's text and text-plus-image accuracy across the five evidence sources, and adds the paired within-question verdict changes across that step. Adding the page image raises accuracy on almost every parser and source combination, the one consistent exception being table evidence on born-digital pages, where three of the four parsers lose ground: a well-parsed table is already a faithful representation, and the image mainly adds an opportunity to misread it. The image also compresses the differences between parsers, with pooled born-digital accuracy spanning 26.6 to 42.3 across parsers on the text channel alone and only 46.0 to 50.2 once the image is added. Regression is otherwise rare, so the image seldom overturns an answer the text channel had already reached.
% GENERATED by tables/scripts/ -- do not edit by hand, re-run the script.
% Numbers read from results/tables/mmlongbench/parser_scan_source.csv.
% Parser fidelity by evidence source and scan status: text source (4) x evidence
% source (5) + a pooled All row per parser.
% Text = TL for the layout parsers, raw embedded text T for PyMuPDF; +Vis adds the page
% image (TLV, or TV for PyMuPDF). The transitions describe that one step.
% Accuracy uses the SAME fixed [0,74] teal scale as findings/parser.tex, so a cell here
% reads against a cell there. The two transition ramps are fitted to the values present:
% R->W (an error) red over [0,10.3], W->R (a recovery) green over [4.9,59.4].
\definecolor{accteal}{HTML}{1BAF7A}
\definecolor{regamber}{HTML}{E23B2E}
\definecolor{recblue}{HTML}{16B84E}
\newcommand{\psacell}[1]{%
  \pgfmathtruncatemacro{\shadeval}{max(0,min(60,(#1-0)/(74-0)*65))}%
  \edef\tmpcol{accteal!\shadeval}%
  \expandafter\cellcolor\expandafter{\tmpcol}#1%
}
\newcommand{\psrcell}[1]{%
  \pgfmathtruncatemacro{\shadeval}{max(0,min(55,(#1-0)/(10.3-0)*55))}%
  \edef\tmpcol{regamber!\shadeval}%
  \expandafter\cellcolor\expandafter{\tmpcol}#1%
}
\newcommand{\pswcell}[1]{%
  \pgfmathtruncatemacro{\shadeval}{max(0,min(60,(#1-4.9)/(59.4-4.9)*65))}%
  \edef\tmpcol{recblue!\shadeval}%
  \expandafter\cellcolor\expandafter{\tmpcol}#1%
}

\begin{table}[ht]
\centering
\small
\setlength{\tabcolsep}{5pt}
\renewcommand{\arraystretch}{1.1}
\adjustbox{max width=\textwidth}{%
\begin{tabular}{@{} ll rrrr @{\hskip 12pt} rrrr @{}}
\toprule
 & & \multicolumn{4}{c}{\textbf{Accuracy (\%)}} & \multicolumn{4}{c}{\textbf{Verdict transition (\%)}} \\
\cmidrule(lr){3-6}\cmidrule(lr){7-10}
 & & \multicolumn{2}{c}{\textbf{Digital}} & \multicolumn{2}{c}{\textbf{Scanned}} & \multicolumn{2}{c}{\textbf{Digital}} & \multicolumn{2}{c}{\textbf{Scanned}} \\
\cmidrule(lr){3-4}\cmidrule(lr){5-6}\cmidrule(lr){7-8}\cmidrule(lr){9-10}
\textbf{Parser} & \textbf{Source} & \textbf{Text} & \textbf{$+$Vis} & \textbf{Text} & \textbf{$+$Vis} & \textbf{R$\rightarrow$W} & \textbf{W$\rightarrow$R} & \textbf{R$\rightarrow$W} & \textbf{W$\rightarrow$R} \\
\midrule
\multirow{6}{*}{PaddleOCR-VL}
 & Chart            & \psacell{23.8}           & \psacell{40.8}           & \psacell{4.2}            & \psacell{54.2}           & \psrcell{3.8}            & \pswcell{20.8}           & \psrcell{0.0}            & \pswcell{50.0} \\
 & Figure           & \psacell{17.4}           & \psacell{37.0}           & \psacell{31.1}           & \psacell{57.5}           & \psrcell{1.6}            & \pswcell{21.2}           & \psrcell{0.9}            & \pswcell{27.4} \\
 & Generalized-text & \psacell{43.8}           & \psacell{47.5}           & \psacell{26.3}           & \psacell{57.9}           & \psrcell{3.8}            & \pswcell{7.5}            & \psrcell{0.0}            & \pswcell{31.6} \\
 & Pure-text        & \psacell{50.9}           & \psacell{54.9}           & \psacell{35.4}           & \psacell{55.4}           & \psrcell{3.5}            & \pswcell{7.5}            & \psrcell{0.0}            & \pswcell{20.0} \\
 & Table            & \psacell{47.6}           & \psacell{46.5}           & \psacell{59.4}           & \psacell{71.9}           & \psrcell{5.9}            & \pswcell{4.9}            & \psrcell{3.1}            & \pswcell{15.6} \\
 & \textbf{All}     & \psacell{40.1}           & \psacell{49.7}           & \psacell{32.8}           & \psacell{61.2}           & \psrcell{3.7}            & \pswcell{13.3}           & \psrcell{1.0}            & \pswcell{29.4} \\
\cmidrule(lr){1-10}
\multirow{6}{*}{MinerU 2.5}
 & Chart            & \psacell{17.7}           & \psacell{38.5}           & \psacell{14.6}           & \psacell{47.9}           & \psrcell{3.1}            & \pswcell{23.8}           & \psrcell{2.1}            & \pswcell{35.4} \\
 & Figure           & \psacell{12.5}           & \psacell{38.6}           & \psacell{25.5}           & \psacell{50.9}           & \psrcell{1.6}            & \pswcell{27.7}           & \psrcell{1.9}            & \pswcell{27.4} \\
 & Generalized-text & \psacell{20.0}           & \psacell{43.8}           & \psacell{18.4}           & \psacell{50.0}           & \psrcell{2.5}            & \pswcell{26.2}           & \psrcell{2.6}            & \pswcell{34.2} \\
 & Pure-text        & \psacell{30.1}           & \psacell{49.1}           & \psacell{30.8}           & \psacell{53.8}           & \psrcell{3.1}            & \pswcell{22.1}           & \psrcell{0.0}            & \pswcell{23.1} \\
 & Table            & \psacell{36.8}           & \psacell{42.7}           & \psacell{56.2}           & \psacell{65.6}           & \psrcell{5.9}            & \pswcell{11.9}           & \psrcell{6.2}            & \pswcell{15.6} \\
 & \textbf{All}     & \psacell{26.6}           & \psacell{46.0}           & \psacell{27.4}           & \psacell{55.7}           & \psrcell{3.4}            & \pswcell{22.8}           & \psrcell{2.0}            & \pswcell{30.3} \\
\cmidrule(lr){1-10}
\multirow{6}{*}{Unlimited-OCR}
 & Chart            & \psacell{23.8}           & \psacell{42.3}           & \psacell{2.1}            & \psacell{58.3}           & \psrcell{3.8}            & \pswcell{22.3}           & \psrcell{0.0}            & \pswcell{56.2} \\
 & Figure           & \psacell{20.7}           & \psacell{38.0}           & \psacell{34.0}           & \psacell{58.5}           & \psrcell{2.7}            & \pswcell{20.1}           & \psrcell{0.9}            & \pswcell{25.5} \\
 & Generalized-text & \psacell{43.8}           & \psacell{52.5}           & \psacell{36.8}           & \psacell{60.5}           & \psrcell{2.5}            & \pswcell{11.2}           & \psrcell{0.0}            & \pswcell{23.7} \\
 & Pure-text        & \psacell{50.4}           & \psacell{53.1}           & \psacell{36.9}           & \psacell{56.9}           & \psrcell{4.4}            & \pswcell{7.1}            & \psrcell{0.0}            & \pswcell{20.0} \\
 & Table            & \psacell{49.7}           & \psacell{48.1}           & \psacell{59.4}           & \psacell{71.9}           & \psrcell{6.5}            & \pswcell{4.9}            & \psrcell{6.2}            & \pswcell{18.8} \\
 & \textbf{All}     & \psacell{42.3}           & \psacell{50.2}           & \psacell{34.3}           & \psacell{63.7}           & \psrcell{4.6}            & \pswcell{12.5}           & \psrcell{1.5}            & \pswcell{30.8} \\
\cmidrule(lr){1-10}
\multirow{6}{*}{PyMuPDF}
 & Chart            & \psacell{24.6}           & \psacell{37.7}           & \psacell{0.0}            & \psacell{54.2}           & \psrcell{3.8}            & \pswcell{16.9}           & \psrcell{0.0}            & \pswcell{54.2} \\
 & Figure           & \psacell{19.6}           & \psacell{41.3}           & \psacell{7.5}            & \psacell{52.8}           & \psrcell{3.3}            & \pswcell{25.0}           & \psrcell{0.9}            & \pswcell{46.2} \\
 & Generalized-text & \psacell{40.0}           & \psacell{47.5}           & \psacell{2.6}            & \psacell{52.6}           & \psrcell{3.8}            & \pswcell{11.2}           & \psrcell{2.6}            & \pswcell{52.6} \\
 & Pure-text        & \psacell{50.0}           & \psacell{54.4}           & \psacell{3.1}            & \psacell{50.8}           & \psrcell{4.4}            & \pswcell{8.8}            & \psrcell{0.0}            & \pswcell{47.7} \\
 & Table            & \psacell{48.1}           & \psacell{44.9}           & \psacell{0.0}            & \psacell{59.4}           & \psrcell{10.3}           & \pswcell{7.0}            & \psrcell{0.0}            & \pswcell{59.4} \\
 & \textbf{All}     & \psacell{40.6}           & \psacell{49.7}           & \psacell{4.0}            & \psacell{57.7}           & \psrcell{4.8}            & \pswcell{13.9}           & \psrcell{0.5}            & \pswcell{54.2} \\
\bottomrule
\end{tabular}%
}
\caption{\small Parser fidelity by evidence source and scan status. Teal deepens with accuracy and green with the vision-step recovery rate; red deepens with the vision-step regression rate. Adding the page image lifts every parser on every source, and it lifts them furthest exactly where the text channel had least to give.}
\label{tab:parser_source_scan}
\end{table}

% {\footnotesize\raggedright
% Accuracy (\%) on oracle pages, Qwen3-VL-8B bf16, medium resolution, answerable pool
% (digital $n{=}646$, scanned $n{=}201$). \textbf{Text} is the text channel alone,
% \textbf{$+$Vis} the same text with the page image added.
% \textbf{R$\rightarrow$W} and \textbf{W$\rightarrow$R} are the paired within-question
% verdict changes across that step: a question answered right on Text and wrong on
% $+$Vis, and the reverse. They are percentages of the cell's own paired questions
% (digital $n{=}80$--$226$ per source, scanned $n{=}32$--$106$), so the scanned
% per-source cells are thin and should be read for direction rather than for level.
% Parser rows partition the pool. Evidence-source rows within a parser overlap, since a
% question citing several sources is counted under each, so their $n$ do not sum and the
% \textbf{All} row pools each question once instead of summing the column.
% \textbf{PyMuPDF on scanned pages sits at the token floor}: raw embedded text is
% near-absent on a scan, so its text-only accuracy reads as ``no usable text reached the
% reasoner'' rather than as a genuine text reading, and its 54.2\% pooled scanned
% recovery rate is the page image supplying the whole answer.\par}

\subsection{Retrieval Across Depth}
\label{app:retrieval_depth}
Table~\ref{tab:retrieval_depth} reports page-level precision, recall, and F1 for all six retrievers at five depths on both corpora, which is the complete sweep behind the operating-point comparison in the main text. Vision retrieval leads text retrieval at every depth on both datasets, and ColQwen3 is the strongest arm on each. F1 peaks at $k{=}1$ for every retriever, since precision falls faster than recall climbs, so a system optimising a single retrieval score should read shallowly while one optimising evidence coverage should read deeply. The gap between the modalities is not an artefact of a weak text arm, since Qwen3-Embedding leads the text arm at every depth beyond the first and still trails ColQwen3 by 22.9 recall points at $k{=}10$ on MMLongBench-Doc. The two corpora are separate pools, so the table is read for the ordering of retrievers within a corpus and for the direction of the depth trend across both.
% GENERATED by tables/scripts/ -- do not edit by hand, re-run the script.
% Numbers read from results/tables/mmlongbench/retrieval_depth_sweep.csv and
% results/tables/longdocurl/ldu_retrieval_overall.csv. The LongDocURL report writes
% fractions where the MMLongBench-Doc one writes percentages, so it is multiplied out.
% Page-level P/R/F1 by retriever x depth, six retrievers on each corpus.
% The joint arms are not shown: a joint arm at depth k returns up to 2k pages, so its
% precision has a different denominator from the single retrievers beside it.
% Amber deepens with the value, on a per-metric scale fitted across both corpora so a
% cell on one corpus reads against a cell on the other:
% P over [9.5,73.3], R over [22.2,88.1], F1 over [15.4,59.5].
% Bold marks each retriever's best-F1 depth within its own corpus.
\definecolor{retamber}{HTML}{C8891B}
\newcommand{\rdpcell}[1]{%
  \pgfmathtruncatemacro{\shadeval}{max(0,min(60,(#1-9.5)/(73.3-9.5)*65))}%
  \edef\tmpcol{retamber!\shadeval}%
  \expandafter\cellcolor\expandafter{\tmpcol}#1%
}
\newcommand{\rdrcell}[1]{%
  \pgfmathtruncatemacro{\shadeval}{max(0,min(60,(#1-22.2)/(88.1-22.2)*65))}%
  \edef\tmpcol{retamber!\shadeval}%
  \expandafter\cellcolor\expandafter{\tmpcol}#1%
}
\newcommand{\rdfshade}[1]{%
  \pgfmathtruncatemacro{\shadeval}{max(0,min(60,(#1-15.4)/(59.5-15.4)*65))}%
  \edef\tmpcol{retamber!\shadeval}%
  \expandafter\cellcolor\expandafter{\tmpcol}%
}
\newcommand{\rdfcell}[1]{\rdfshade{#1}#1}
\newcommand{\rdfbest}[1]{\rdfshade{#1}\textbf{#1}}

\begin{table}[ht]
\centering
\small
\setlength{\tabcolsep}{5pt}
\renewcommand{\arraystretch}{1.1}
\adjustbox{max width=\textwidth}{%
\begin{tabular}{l@{\hskip 14pt}rrrrr@{\hskip 14pt}rrrrr@{\hskip 14pt}rrrrr}
\toprule
& \multicolumn{5}{c}{\textbf{Precision}}
& \multicolumn{5}{c}{\textbf{Recall}}
& \multicolumn{5}{c}{\textbf{F1}} \\
\cmidrule(lr){2-6}\cmidrule(lr){7-11}\cmidrule(lr){12-16}
\textbf{Retriever}
& \textbf{1} & \textbf{3} & \textbf{5} & \textbf{7} & \textbf{10}
& \textbf{1} & \textbf{3} & \textbf{5} & \textbf{7} & \textbf{10}
& \textbf{1} & \textbf{3} & \textbf{5} & \textbf{7} & \textbf{10} \\
\midrule
\multicolumn{16}{@{}l}{\textbf{MMLongBench-Doc} ($n{=}847$)} \\
\multicolumn{16}{@{}l}{\quad\emph{Text}} \\
\quad\quad BM25            & \rdpcell{29.5}   & \rdpcell{18.5}   & \rdpcell{14.0}   & \rdpcell{11.4}   & \rdpcell{9.5}    & \rdrcell{22.2}   & \rdrcell{37.8}   & \rdrcell{46.0}   & \rdrcell{50.8}   & \rdrcell{58.5}   & \rdfbest{24.2}   & \rdfcell{23.4}   & \rdfcell{20.1}   & \rdfcell{17.5}   & \rdfcell{15.4} \\
\quad\quad BGE-M3          & \rdpcell{33.8}   & \rdpcell{18.7}   & \rdpcell{14.7}   & \rdpcell{12.0}   & \rdpcell{10.0}   & \rdrcell{25.5}   & \rdrcell{38.8}   & \rdrcell{48.2}   & \rdrcell{53.2}   & \rdrcell{61.1}   & \rdfbest{27.9}   & \rdfcell{23.9}   & \rdfcell{21.2}   & \rdfcell{18.4}   & \rdfcell{16.3} \\
\quad\quad Qwen3-Embedding & \rdpcell{33.5}   & \rdpcell{19.9}   & \rdpcell{15.0}   & \rdpcell{12.6}   & \rdpcell{10.5}   & \rdrcell{25.5}   & \rdrcell{41.3}   & \rdrcell{49.4}   & \rdrcell{56.0}   & \rdrcell{64.5}   & \rdfbest{27.9}   & \rdfcell{25.4}   & \rdfcell{21.7}   & \rdfcell{19.4}   & \rdfcell{17.1} \\
\multicolumn{16}{@{}l}{\quad\emph{Vision}} \\
\quad\quad ColModernVBERT  & \rdpcell{58.8}   & \rdpcell{31.5}   & \rdpcell{22.7}   & \rdpcell{17.9}   & \rdpcell{14.0}   & \rdrcell{46.0}   & \rdrcell{64.6}   & \rdrcell{72.5}   & \rdrcell{77.2}   & \rdrcell{81.7}   & \rdfbest{49.6}   & \rdfcell{39.6}   & \rdfcell{32.2}   & \rdfcell{27.1}   & \rdfcell{22.2} \\
\quad\quad ColQwen2.5      & \rdpcell{63.3}   & \rdpcell{33.8}   & \rdpcell{23.6}   & \rdpcell{18.5}   & \rdpcell{14.4}   & \rdrcell{48.8}   & \rdrcell{68.3}   & \rdrcell{74.9}   & \rdrcell{79.2}   & \rdrcell{83.8}   & \rdfbest{52.8}   & \rdfcell{42.5}   & \rdfcell{33.5}   & \rdfcell{27.9}   & \rdfcell{22.9} \\
\quad\quad ColQwen3        & \rdpcell{70.2}   & \rdpcell{36.8}   & \rdpcell{25.5}   & \rdpcell{19.9}   & \rdpcell{15.2}   & \rdrcell{54.1}   & \rdrcell{73.9}   & \rdrcell{80.6}   & \rdrcell{83.9}   & \rdrcell{87.4}   & \rdfbest{58.6}   & \rdfcell{46.0}   & \rdfcell{36.1}   & \rdfcell{29.9}   & \rdfcell{24.1} \\
\cmidrule(lr){1-16}
\multicolumn{16}{@{}l}{\textbf{LongDocURL} ($n{=}2{,}317$)} \\
\multicolumn{16}{@{}l}{\quad\emph{Text}} \\
\quad\quad BM25            & \rdpcell{49.6}   & \rdpcell{27.0}   & \rdpcell{19.0}   & \rdpcell{14.8}   & \rdpcell{11.3}   & \rdrcell{35.6}   & \rdrcell{53.9}   & \rdrcell{61.1}   & \rdrcell{65.6}   & \rdrcell{69.5}   & \rdfbest{39.7}   & \rdfcell{34.3}   & \rdfcell{27.7}   & \rdfcell{23.3}   & \rdfcell{18.7} \\
\quad\quad BGE-M3          & \rdpcell{47.6}   & \rdpcell{27.8}   & \rdpcell{19.8}   & \rdpcell{15.5}   & \rdpcell{11.9}   & \rdrcell{33.4}   & \rdrcell{54.2}   & \rdrcell{62.3}   & \rdrcell{67.3}   & \rdrcell{72.6}   & \rdfbest{37.6}   & \rdfcell{35.0}   & \rdfcell{28.7}   & \rdfcell{24.1}   & \rdfcell{19.8} \\
\quad\quad Qwen3-Embedding & \rdpcell{48.9}   & \rdpcell{28.8}   & \rdpcell{20.6}   & \rdpcell{16.4}   & \rdpcell{12.5}   & \rdrcell{34.1}   & \rdrcell{56.0}   & \rdrcell{64.8}   & \rdrcell{70.6}   & \rdrcell{75.5}   & \rdfbest{38.5}   & \rdfcell{36.3}   & \rdfcell{29.9}   & \rdfcell{25.5}   & \rdfcell{20.7} \\
\multicolumn{16}{@{}l}{\quad\emph{Vision}} \\
\quad\quad ColModernVBERT  & \rdpcell{64.7}   & \rdpcell{33.5}   & \rdpcell{23.2}   & \rdpcell{17.9}   & \rdpcell{13.6}   & \rdrcell{47.5}   & \rdrcell{67.3}   & \rdrcell{74.5}   & \rdrcell{78.6}   & \rdrcell{83.1}   & \rdfbest{52.7}   & \rdfcell{42.7}   & \rdfcell{33.9}   & \rdfcell{28.0}   & \rdfcell{22.5} \\
\quad\quad ColQwen2.5      & \rdpcell{68.1}   & \rdpcell{35.8}   & \rdpcell{24.8}   & \rdpcell{18.9}   & \rdpcell{14.2}   & \rdrcell{49.3}   & \rdrcell{70.6}   & \rdrcell{78.0}   & \rdrcell{81.5}   & \rdrcell{85.2}   & \rdfbest{54.9}   & \rdfcell{45.2}   & \rdfcell{35.9}   & \rdfcell{29.4}   & \rdfcell{23.4} \\
\quad\quad ColQwen3        & \rdpcell{73.3}   & \rdpcell{38.2}   & \rdpcell{25.9}   & \rdpcell{19.8}   & \rdpcell{14.8}   & \rdrcell{53.6}   & \rdrcell{74.5}   & \rdrcell{81.1}   & \rdrcell{84.6}   & \rdrcell{88.1}   & \rdfbest{59.5}   & \rdfcell{48.1}   & \rdfcell{37.5}   & \rdfcell{30.7}   & \rdfcell{24.4} \\
\bottomrule
\end{tabular}%
}
\caption{\small Page-level retrieval across depth on both corpora. Amber deepens with the value on a per-metric scale and bold marks each retriever's best-F1 depth. Vision retrieval leads text at every depth on both corpora, and recall climbs with $k$ as precision falls.}
\label{tab:retrieval_depth}
\end{table}

% {\footnotesize\raggedright
% Page-level precision, recall and F1 (\%) at retrieval depths $k \in \{1,3,5,7,10\}$,
% rendered at 200 DPI on each corpus's answerable pool. Every cell is computed from the
% retriever's stored full page ranking, so the sweep is complete rather than interpolated.
% Each metric block is shaded on its own scale, fitted across both corpora, so amber
% deepens with the value within that metric and the depth trend is visible per block:
% precision falls with depth as recall climbs. Bold marks each retriever's best-F1 depth
% within its own corpus.
% \textbf{The joint arms are not shown.} A joint arm at depth $k$ returns up to $2k$
% deduplicated pages, so its precision is divided by roughly twice the denominator of the
% single retrievers beside it and the column stops being one comparison; the pooled
% retrieval table reports them where that trade is the point.
% \textbf{The two corpora are not one pool} and their absolute levels are not
% interchangeable: LongDocURL's documents are longer, so a page drawn at random is a worse
% guess there, and its precision is nonetheless higher because more of its questions cite
% a single page. Read the ordering of retrievers within a corpus and the direction of the
% depth trend across both.
% \textbf{Vision retrieval leads text at every depth on both corpora} and ColQwen3 is
% best on F1 throughout. The gap is not an artefact of a weak text arm: Qwen3-Embedding is
% the strongest text retriever at nearly every fixed depth and still trails ColQwen3 by a
% wide margin.\par}

\subsection{Prompting by Evidence Source}
\label{app:prompting_source}
Table~\ref{tab:prompting_source} opens the prompting sweep up by evidence source, reporting answerable accuracy, abstention on the answerable pool, and abstention on the unanswerable pool for six prompt modes at four rungs. The abstention instruction is the only prompt that materially raises correct refusal, lifting unanswerable abstention from roughly 13\% uninstructed to between 75 and 82\% across the four rungs, but it does so indiscriminately: abstention on the answerable pool rises with it at every rung, and the per-source rows show that the coupling is not confined to the harder evidence types. The two chain-of-thought modes leave both abstention rates near the uninstructed baseline, and \texttt{cot} is the only mode that matches \texttt{none} on answerable accuracy at TLV. Evidence-source ordering is stable across prompt modes, with pure-text and table questions strongest and chart and figure questions weakest, and adding the page image closes most of that spread. The unanswerable block carries no per-source rows, an unanswerable question having no gold evidence pages from which a source could be read, and it is scored under the original judge rather than the error-taxonomy judge used for the two answerable blocks, so levels are compared within a block rather than across the table.
% GENERATED by tables/scripts/ -- do not edit by hand, re-run the script.
% Numbers read from results/tables/analysis/faithfulness_answerable_accuracy_source.csv
% (the two answerable blocks) and results/tables/mmlongbench/
% faithfulness_unanswerable_abstention.csv (the third block).
% Prompt mode (6) x evidence source (5) + a pooled All row per mode, four rungs each.
% Shade ranges are fitted to the values present:
% teal marks the good direction -- answerable accuracy over [6.9,54.2] and correct
% unanswerable abstention over [5.3,81.6]; warm marks answerable abstention over
% [0,76.2], which is an error wherever the rung carried the evidence.
% The unanswerable block has no per-source rows and cannot have any: an unanswerable
% question has no gold evidence pages, so the pool cannot be split by source at all.
\definecolor{accteal}{HTML}{1BAF7A}
\definecolor{recwarm}{HTML}{D95926}
\newcommand{\pmacell}[1]{%
  \pgfmathtruncatemacro{\shadeval}{max(0,min(60,(#1-6.9)/(54.2-6.9)*65))}%
  \edef\tmpcol{accteal!\shadeval}%
  \expandafter\cellcolor\expandafter{\tmpcol}#1%
}
\newcommand{\pmfcell}[1]{%
  \pgfmathtruncatemacro{\shadeval}{max(0,min(55,(#1-0)/(76.2-0)*55))}%
  \edef\tmpcol{recwarm!\shadeval}%
  \expandafter\cellcolor\expandafter{\tmpcol}#1%
}
\newcommand{\pmucell}[1]{%
  \pgfmathtruncatemacro{\shadeval}{max(0,min(60,(#1-5.3)/(81.6-5.3)*65))}%
  \edef\tmpcol{accteal!\shadeval}%
  \expandafter\cellcolor\expandafter{\tmpcol}#1%
}

\begin{table}[ht]
\centering
\small
\setlength{\tabcolsep}{5pt}
\renewcommand{\arraystretch}{1.05}
\adjustbox{max width=\textwidth}{%
\begin{tabular}{ll@{\hskip 14pt}rrrr@{\hskip 18pt}rrrr@{\hskip 18pt}rrrr}
\toprule
& & \multicolumn{4}{c}{\textbf{Answerable acc.}}
& \multicolumn{4}{c}{\textbf{Abstention (ans.)}}
& \multicolumn{4}{c}{\textbf{Abstention (unans.)}} \\
\cmidrule(lr){3-6}\cmidrule(lr){7-10}\cmidrule(lr){11-14}
\textbf{Prompt} & \textbf{Source}
& \textbf{T} & \textbf{TL} & \textbf{TLV} & \textbf{V}
& \textbf{T} & \textbf{TL} & \textbf{TLV} & \textbf{V}
& \textbf{T} & \textbf{TL} & \textbf{TLV} & \textbf{V} \\
\midrule
\multirow{6}{*}{none}
 & Chart            & \pmacell{16.3}         & \pmacell{16.9}         & \pmacell{42.7}         & \pmacell{40.4}         & \pmfcell{32.6}         & \pmfcell{28.1}         & \pmfcell{7.3}          & \pmfcell{6.2}          &                        &                        &                        &  \\
 & Figure           & \pmacell{14.5}         & \pmacell{22.1}         & \pmacell{42.4}         & \pmacell{39.0}         & \pmfcell{36.6}         & \pmfcell{24.8}         & \pmfcell{7.6}          & \pmfcell{7.9}          &                        &                        &                        &  \\
 & Generalized-text & \pmacell{26.3}         & \pmacell{33.9}         & \pmacell{49.2}         & \pmacell{43.2}         & \pmfcell{20.3}         & \pmfcell{10.2}         & \pmfcell{0.8}          & \pmfcell{4.2}          &                        &                        &                        &  \\
 & Pure-text        & \pmacell{36.9}         & \pmacell{44.4}         & \pmacell{51.0}         & \pmacell{42.3}         & \pmfcell{23.0}         & \pmfcell{11.5}         & \pmfcell{7.3}          & \pmfcell{8.0}          &                        &                        &                        &  \\
 & Table            & \pmacell{35.6}         & \pmacell{44.7}         & \pmacell{44.7}         & \pmacell{34.1}         & \pmfcell{15.9}         & \pmfcell{7.2}          & \pmfcell{6.7}          & \pmfcell{9.1}          &                        &                        &                        &  \\
 & \textbf{All}     & \pmacell{29.0}         & \pmacell{35.7}         & \pmacell{49.5}         & \pmacell{42.6}         & \pmfcell{26.3}         & \pmfcell{17.2}         & \pmfcell{6.2}          & \pmfcell{7.1}          & \pmucell{13.1}         & \pmucell{13.1}         & \pmucell{14.3}         & \pmucell{13.5} \\
\cmidrule(lr){1-14}
\multirow{6}{*}{grounded}
 & Chart            & \pmacell{15.2}         & \pmacell{14.6}         & \pmacell{31.5}         & \pmacell{28.1}         & \pmfcell{28.1}         & \pmfcell{23.0}         & \pmfcell{5.1}          & \pmfcell{5.6}          &                        &                        &                        &  \\
 & Figure           & \pmacell{11.4}         & \pmacell{21.4}         & \pmacell{40.7}         & \pmacell{37.9}         & \pmfcell{34.1}         & \pmfcell{21.4}         & \pmfcell{7.2}          & \pmfcell{7.9}          &                        &                        &                        &  \\
 & Generalized-text & \pmacell{19.5}         & \pmacell{29.7}         & \pmacell{43.2}         & \pmacell{44.9}         & \pmfcell{17.8}         & \pmfcell{6.8}          & \pmfcell{0.0}          & \pmfcell{0.8}          &                        &                        &                        &  \\
 & Pure-text        & \pmacell{35.2}         & \pmacell{42.3}         & \pmacell{49.0}         & \pmacell{40.6}         & \pmfcell{19.9}         & \pmfcell{10.1}         & \pmfcell{7.0}          & \pmfcell{6.3}          &                        &                        &                        &  \\
 & Table            & \pmacell{27.9}         & \pmacell{37.5}         & \pmacell{38.5}         & \pmacell{28.4}         & \pmfcell{17.3}         & \pmfcell{13.0}         & \pmfcell{7.2}          & \pmfcell{8.2}          &                        &                        &                        &  \\
 & \textbf{All}     & \pmacell{24.9}         & \pmacell{32.3}         & \pmacell{43.7}         & \pmacell{37.9}         & \pmfcell{24.1}         & \pmfcell{16.2}         & \pmfcell{6.0}          & \pmfcell{6.5}          & \pmucell{22.1}         & \pmucell{6.6}          & \pmucell{8.2}          & \pmucell{5.3} \\
\cmidrule(lr){1-14}
\multirow{6}{*}{abstain}
 & Chart            & \pmacell{13.5}         & \pmacell{11.8}         & \pmacell{26.4}         & \pmacell{22.5}         & \pmfcell{52.2}         & \pmfcell{59.0}         & \pmfcell{25.3}         & \pmfcell{27.5}         &                        &                        &                        &  \\
 & Figure           & \pmacell{6.9}          & \pmacell{13.8}         & \pmacell{34.1}         & \pmacell{30.3}         & \pmfcell{76.2}         & \pmfcell{65.2}         & \pmfcell{35.5}         & \pmfcell{38.3}         &                        &                        &                        &  \\
 & Generalized-text & \pmacell{19.5}         & \pmacell{25.4}         & \pmacell{41.5}         & \pmacell{37.3}         & \pmfcell{55.9}         & \pmfcell{44.1}         & \pmfcell{19.5}         & \pmfcell{25.4}         &                        &                        &                        &  \\
 & Pure-text        & \pmacell{31.7}         & \pmacell{37.8}         & \pmacell{43.7}         & \pmacell{34.3}         & \pmfcell{42.5}         & \pmfcell{36.0}         & \pmfcell{24.1}         & \pmfcell{31.1}         &                        &                        &                        &  \\
 & Table            & \pmacell{24.5}         & \pmacell{37.0}         & \pmacell{38.5}         & \pmacell{26.4}         & \pmfcell{43.3}         & \pmfcell{37.0}         & \pmfcell{30.3}         & \pmfcell{38.9}         &                        &                        &                        &  \\
 & \textbf{All}     & \pmacell{21.7}         & \pmacell{27.9}         & \pmacell{40.0}         & \pmacell{31.9}         & \pmfcell{52.8}         & \pmfcell{47.7}         & \pmfcell{27.7}         & \pmfcell{32.3}         & \pmucell{81.1}         & \pmucell{79.9}         & \pmucell{74.6}         & \pmucell{81.6} \\
\cmidrule(lr){1-14}
\multirow{6}{*}{abstain\_bal.}
 & Chart            & \pmacell{12.4}         & \pmacell{11.8}         & \pmacell{28.1}         & \pmacell{23.0}         & \pmfcell{44.9}         & \pmfcell{51.7}         & \pmfcell{18.5}         & \pmfcell{24.2}         &                        &                        &                        &  \\
 & Figure           & \pmacell{6.9}          & \pmacell{14.5}         & \pmacell{35.5}         & \pmacell{32.8}         & \pmfcell{71.7}         & \pmfcell{62.4}         & \pmfcell{31.4}         & \pmfcell{33.4}         &                        &                        &                        &  \\
 & Generalized-text & \pmacell{19.5}         & \pmacell{28.8}         & \pmacell{41.5}         & \pmacell{39.8}         & \pmfcell{52.5}         & \pmfcell{41.5}         & \pmfcell{18.6}         & \pmfcell{21.2}         &                        &                        &                        &  \\
 & Pure-text        & \pmacell{31.0}         & \pmacell{37.8}         & \pmacell{45.5}         & \pmacell{34.3}         & \pmfcell{39.7}         & \pmfcell{33.2}         & \pmfcell{19.2}         & \pmfcell{25.5}         &                        &                        &                        &  \\
 & Table            & \pmacell{24.0}         & \pmacell{36.1}         & \pmacell{35.6}         & \pmacell{28.4}         & \pmfcell{38.5}         & \pmfcell{30.3}         & \pmfcell{25.5}         & \pmfcell{34.6}         &                        &                        &                        &  \\
 & \textbf{All}     & \pmacell{21.7}         & \pmacell{28.6}         & \pmacell{40.9}         & \pmacell{33.5}         & \pmfcell{48.2}         & \pmfcell{43.5}         & \pmfcell{23.2}         & \pmfcell{27.9}         & \pmucell{75.4}         & \pmucell{77.0}         & \pmucell{68.9}         & \pmucell{76.6} \\
\cmidrule(lr){1-14}
\multirow{6}{*}{cot}
 & Chart            & \pmacell{19.7}         & \pmacell{16.3}         & \pmacell{43.3}         & \pmacell{38.2}         & \pmfcell{16.3}         & \pmfcell{42.1}         & \pmfcell{9.6}          & \pmfcell{5.6}          &                        &                        &                        &  \\
 & Figure           & \pmacell{11.0}         & \pmacell{17.9}         & \pmacell{43.8}         & \pmacell{43.8}         & \pmfcell{15.9}         & \pmfcell{26.2}         & \pmfcell{6.9}          & \pmfcell{5.9}          &                        &                        &                        &  \\
 & Generalized-text & \pmacell{25.4}         & \pmacell{34.7}         & \pmacell{49.2}         & \pmacell{54.2}         & \pmfcell{3.4}          & \pmfcell{11.0}         & \pmfcell{3.4}          & \pmfcell{3.4}          &                        &                        &                        &  \\
 & Pure-text        & \pmacell{32.4}         & \pmacell{42.3}         & \pmacell{51.0}         & \pmacell{46.2}         & \pmfcell{13.9}         & \pmfcell{15.7}         & \pmfcell{7.7}          & \pmfcell{5.6}          &                        &                        &                        &  \\
 & Table            & \pmacell{37.5}         & \pmacell{48.6}         & \pmacell{51.4}         & \pmacell{44.0}         & \pmfcell{10.6}         & \pmfcell{12.0}         & \pmfcell{7.2}          & \pmfcell{5.3}          &                        &                        &                        &  \\
 & \textbf{All}     & \pmacell{26.9}         & \pmacell{34.1}         & \pmacell{50.2}         & \pmacell{46.9}         & \pmfcell{13.3}         & \pmfcell{21.7}         & \pmfcell{6.6}          & \pmfcell{5.3}          & \pmucell{8.6}          & \pmucell{13.9}         & \pmucell{11.9}         & \pmucell{5.7} \\
\cmidrule(lr){1-14}
\multirow{6}{*}{extract\_cot}
 & Chart            & \pmacell{13.1}         & \pmacell{16.3}         & \pmacell{36.6}         & \pmacell{35.6}         & \pmfcell{19.3}         & \pmfcell{37.6}         & \pmfcell{6.3}          & \pmfcell{7.3}          &                        &                        &                        &  \\
 & Figure           & \pmacell{10.3}         & \pmacell{20.1}         & \pmacell{43.3}         & \pmacell{39.4}         & \pmfcell{27.6}         & \pmfcell{36.5}         & \pmfcell{11.1}         & \pmfcell{11.8}         &                        &                        &                        &  \\
 & Generalized-text & \pmacell{21.2}         & \pmacell{34.2}         & \pmacell{51.3}         & \pmacell{51.7}         & \pmfcell{7.6}          & \pmfcell{18.8}         & \pmfcell{2.6}          & \pmfcell{4.2}          &                        &                        &                        &  \\
 & Pure-text        & \pmacell{34.5}         & \pmacell{40.4}         & \pmacell{48.2}         & \pmacell{41.2}         & \pmfcell{17.1}         & \pmfcell{19.3}         & \pmfcell{9.9}          & \pmfcell{11.6}         &                        &                        &                        &  \\
 & Table            & \pmacell{37.7}         & \pmacell{46.2}         & \pmacell{49.0}         & \pmacell{36.9}         & \pmfcell{13.0}         & \pmfcell{14.4}         & \pmfcell{9.6}          & \pmfcell{17.0}         &                        &                        &                        &  \\
 & \textbf{All}     & \pmacell{26.1}         & \pmacell{33.5}         & \pmacell{47.5}         & \pmacell{41.6}         & \pmfcell{19.4}         & \pmfcell{26.5}         & \pmfcell{8.8}          & \pmfcell{11.6}         & \pmucell{8.6}          & \pmucell{18.0}         & \pmucell{18.4}         & \pmucell{18.4} \\
\bottomrule
\end{tabular}%
}
\caption{\small Faithfulness across six prompt modes by evidence source. Teal deepens with the good-direction rate and warm with abstention on the answerable pool; the abstention instruction couples correct refusal to false refusal on every evidence source, and no prompt separates the two.}
\label{tab:prompting_source}
\end{table}

% {\footnotesize\raggedright
% Judge-scored rates (\%) across six prompt modes, five evidence sources and four rungs.
% Answerable accuracy and answerable abstention run on oracle pages over the answerable
% pool; unanswerable abstention runs on the bm25 $k{=}3$ page set over the unanswerable
% pool ($n{=}244$ per rung), since an unanswerable question has no gold page to select.
% \textbf{The two answerable blocks are judged under the error-taxonomy prompt}, the one
% that defines the abstained verdict, and are directly comparable with
% Table~\ref{tab:prompting}. \textbf{The unanswerable block has no re-judged counterpart
% and is the original judge's}, so compare within a block and not across the rule.
% Per-mode \textbf{All} rows pool every question once; source rows overlap, since a
% question citing several sources appears under each, so their $n$ do not sum to the
% \textbf{All} row above them.
% \textbf{The unanswerable block has no per-source rows because the split does not
% exist}: an unanswerable question carries no gold evidence pages, so there is nothing to
% read a source off, and only the pooled row of that block is defined.
% \textbf{Answerable abstention is not simply an error rate.} The oracle pages carry the
% evidence but the rung need not: \textbf{T} is the embedded text layer and a scanned
% page has none, so at \textbf{T} on a scanned document declining is the correct response
% to what the rung delivered. That is why the column is shaded for magnitude rather than as
% a fault, and why the born-digital component in the source report is the false-abstention
% rate proper.\par}

\subsection{Reasoner Scale by Evidence-Page Count}
\label{app:reasoner_hop}
Table~\ref{tab:reasoner_hop} reports accuracy by the number of gold evidence pages a question cites, for four Qwen3-VL sizes at the two rungs that carry the page image. Scale buys multi-page accuracy steadily, and buys most where the image carries the evidence, with the 32B gaining 7.8 points over the 8B at TLV and 12.5 at V. What scale does not buy is integration, since single-page accuracy rises just as fast and the deficit between multi-page and single-page questions is essentially unchanged across the size ladder. Read alongside Table~\ref{tab:reasoner}, this supports treating cross-page integration as a distinct failure mode rather than a capability that additional parameters resolve. The 3+ bucket is thin at every size and is read for trend rather than for level.
% GENERATED by tables/scripts/ -- do not edit by hand, re-run the script.
% Numbers read from results/tables/mmlongbench/reasoner_hop.csv.
% Accuracy by gold evidence-page bucket, per Qwen3-VL size, at TLV and V.
% T and TL are not shown: both are text-only, and the question here is whether
% parameters buy multi-page accuracy on the rungs that carry the page image.
% Accuracy takes the SAME fixed [14,74] teal scale as findings/ladder.tex and
% discussion/reasoner.tex, so a cell here reads against a cell there.
% The last column of each rung is NOT an accuracy: it is multi-page accuracy read
% against the 8B at the same rung, so it carries its own diverging pair over
% [0,12.8] points -- green above the 8B, warm below it.
\definecolor{accteal}{HTML}{1BAF7A}
\definecolor{recwarm}{HTML}{D95926}
\definecolor{gainmoss}{HTML}{16805F}
\newcommand{\rhacell}[1]{%
  \pgfmathtruncatemacro{\shadeval}{max(0,min(60,(#1-14)/(74-14)*65))}%
  \edef\tmpcol{accteal!\shadeval}%
  \expandafter\cellcolor\expandafter{\tmpcol}#1%
}
\newcommand{\rhugain}[1]{%
  \pgfmathtruncatemacro{\shadeval}{max(0,min(55,(#1-0)/(12.8-0)*60))}%
  \edef\tmpcol{gainmoss!\shadeval}%
  \expandafter\cellcolor\expandafter{\tmpcol}$+$#1%
}
\newcommand{\rhdgain}[1]{%
  \pgfmathtruncatemacro{\shadeval}{max(0,min(55,(#1-0)/(12.8-0)*60))}%
  \edef\tmpcol{recwarm!\shadeval}%
  \expandafter\cellcolor\expandafter{\tmpcol}$-$#1%
}

\begin{table}[ht]
\centering
\small
\setlength{\tabcolsep}{4pt}
\renewcommand{\arraystretch}{1.15}
\adjustbox{max width=\textwidth}{%
\begin{tabular}{@{}l rrrrr @{\hskip 14pt} rrrrr@{}}
\toprule
 & \multicolumn{5}{c}{\textbf{TLV}} & \multicolumn{5}{c}{\textbf{V}} \\
\cmidrule(lr){2-6}\cmidrule(lr){7-11}
\textbf{Size}
 & \textbf{1} & \textbf{2} & \textbf{3+} & \textbf{M} & \textbf{$\Delta$M}
 & \textbf{1} & \textbf{2} & \textbf{3+} & \textbf{M} & \textbf{$\Delta$M} \\
\midrule
2B   & \rhacell{48.1}   & \rhacell{27.4}   & \rhacell{25.2}   & \rhacell{26.7}   & \rhdgain{11.9}   & \rhacell{40.7}   & \rhacell{20.7}   & \rhacell{22.5}   & \rhacell{21.3}   & \rhdgain{12.8} \\
4B   & \rhacell{60.5}   & \rhacell{42.3}   & \rhacell{32.4}   & \rhacell{39.2}   & \rhugain{0.6}    & \rhacell{50.4}   & \rhacell{34.9}   & \rhacell{31.5}   & \rhacell{33.8}   & \rhdgain{0.3} \\
8B   & \rhacell{64.6}   & \rhacell{38.6}   & \rhacell{38.7}   & \rhacell{38.6}   & --               & \rhacell{55.7}   & \rhacell{35.3}   & \rhacell{31.5}   & \rhacell{34.1}   & -- \\
32B  & \rhacell{71.5}   & \rhacell{49.0}   & \rhacell{40.9}   & \rhacell{46.4}   & \rhugain{7.8}    & \rhacell{66.0}   & \rhacell{48.5}   & \rhacell{42.3}   & \rhacell{46.6}   & \rhugain{12.5} \\
\bottomrule
\end{tabular}%
}
\caption{\small Accuracy by gold evidence-page count for each reasoner size. \textbf{M} pools every multi-page question and $\Delta$\textbf{M} reads it against the 8B at the same rung. Scale buys multi-page accuracy, most of it on the image encodings, but it raises single-page accuracy just as fast, so the multi-page deficit does not close.}
\label{tab:reasoner_hop}
\end{table}

% {\footnotesize\raggedright
% Accuracy (\%) on oracle pages at medium resolution, by the number of gold evidence
% pages the question cites, for each Qwen3-VL size ($n{=}826$ per cell; $825$ for the
% 32B at \textbf{TLV}, whose 4+ bucket is one cell short). Teal deepens with accuracy on
% the same fixed scale as Table~\ref{tab:ladder_source} and Table~\ref{tab:reasoner}.
% \textbf{M} is cumulative multi-page evidence, two or more gold pages, so it pools the
% \textbf{2} and \textbf{3+} columns rather than being a fourth disjoint bucket.
% $\Delta$\textbf{M} is that column read against the 8B at the same rung, green above it
% and warm below it; it is a point difference rather than an accuracy, which is why it
% carries its own scale, and the 8B row prints a dash rather than a shaded zero.
% \textbf{Scale buys multi-page accuracy and buys most of it from the image}: the 32B
% gains 7.8 points of multi-page accuracy over the 8B at \textbf{TLV} and 12.5 at
% \textbf{V}. \textbf{What it does not buy is integration}, since single-page accuracy
% rises just as fast and the gap between the two is no smaller at 32B than at 2B. The
% \textbf{3+} bucket is thin at every size and should be read for trend rather than for
% level.\par}


% \clearpage
% \section{Extended Error Analysis Tables}
% \label{app:extended_analysis}

% \subsection{Failure Modes by Prompt Mode}
% \label{app:error_profile}
% Table~\ref{tab:error_profile} is the tabular form of Figure~\ref{fig:error_profile}. The figure pools five of the ten categories, equivocal, the two list errors, unit or format mismatch, and other, into a single residual band, while the table prints each of them separately alongside the \emph{All incorrect} row they sum to. The two blocks are not comparable in level, since the answerable pool runs on oracle pages and the unanswerable pool on the BM25 $k{=}3$ page set, an unanswerable question having no gold page to select.
% % GENERATED by tables/scripts/ -- do not edit by hand, re-run the script.
% Numbers read from results/tables/analysis/*.csv.
% error_profile_prompt, error_profile_prompt_unanswerable (results.errcat.jsonl)
% Shade ranges are fitted to the values present:
% Share of the pool, at TLV (\%): recwarm over [0.0, 60.0]
\definecolor{recwarm}{HTML}{D95926}
\begin{table*}[h]
\centering
\small
\setlength{\tabcolsep}{5pt}
\adjustbox{max width=\textwidth}{%
\begin{tabular}{l@{\hskip 14pt}rrrrrr}
\toprule
 & \multicolumn{6}{c}{\textbf{Share of the pool, at TLV (\%)}} \\
\cmidrule(lr){2-7}
\textbf{Error category} & \textbf{none} & \textbf{grounded} & \textbf{abstain} & \textbf{abstain\_bal.} & \textbf{cot} & \textbf{extract\_cot} \\
\midrule
\multicolumn{7}{l}{\emph{answerable pool (oracle pages)}} \\
insufficient evidence & \cellcolor{recwarm!6}6.1 & \cellcolor{recwarm!6}6.0 & \cellcolor{recwarm!28}\textbf{27.6} & \cellcolor{recwarm!23}\textbf{23.0} & \cellcolor{recwarm!7}6.6 & \cellcolor{recwarm!9}8.8 \\
no answer reached & \cellcolor{recwarm!7}6.8 & \cellcolor{recwarm!0}0.0 & \cellcolor{recwarm!0}0.0 & \cellcolor{recwarm!0}0.0 & \cellcolor{recwarm!4}4.4 & \cellcolor{recwarm!3}3.3 \\
equivocal & \cellcolor{recwarm!1}1.1 & \cellcolor{recwarm!0}0.1 & \cellcolor{recwarm!0}0.0 & \cellcolor{recwarm!0}0.0 & \cellcolor{recwarm!0}0.0 & \cellcolor{recwarm!0}0.1 \\
wrong count & \cellcolor{recwarm!14}\textbf{14.4} & \cellcolor{recwarm!19}\textbf{18.8} & \cellcolor{recwarm!11}10.9 & \cellcolor{recwarm!12}11.8 & \cellcolor{recwarm!15}\textbf{14.8} & \cellcolor{recwarm!13}\textbf{12.7} \\
wrong value & \cellcolor{recwarm!9}9.0 & \cellcolor{recwarm!15}14.6 & \cellcolor{recwarm!9}8.8 & \cellcolor{recwarm!11}10.8 & \cellcolor{recwarm!8}8.5 & \cellcolor{recwarm!10}10.1 \\
wrong entity & \cellcolor{recwarm!8}7.7 & \cellcolor{recwarm!10}9.6 & \cellcolor{recwarm!7}7.1 & \cellcolor{recwarm!7}7.3 & \cellcolor{recwarm!7}7.1 & \cellcolor{recwarm!9}8.7 \\
incomplete list & \cellcolor{recwarm!2}1.8 & \cellcolor{recwarm!3}2.6 & \cellcolor{recwarm!2}1.9 & \cellcolor{recwarm!2}2.0 & \cellcolor{recwarm!3}2.9 & \cellcolor{recwarm!2}2.2 \\
over-inclusive list & \cellcolor{recwarm!2}2.2 & \cellcolor{recwarm!3}2.8 & \cellcolor{recwarm!2}2.0 & \cellcolor{recwarm!2}2.5 & \cellcolor{recwarm!3}2.9 & \cellcolor{recwarm!3}3.4 \\
unit / format mismatch & \cellcolor{recwarm!1}1.3 & \cellcolor{recwarm!2}1.6 & \cellcolor{recwarm!1}1.4 & \cellcolor{recwarm!1}1.3 & \cellcolor{recwarm!2}2.3 & \cellcolor{recwarm!2}2.3 \\
other & \cellcolor{recwarm!0}0.1 & \cellcolor{recwarm!0}0.1 & \cellcolor{recwarm!0}0.1 & \cellcolor{recwarm!0}0.1 & \cellcolor{recwarm!0}0.4 & \cellcolor{recwarm!1}0.8 \\
\emph{All incorrect} & \cellcolor{recwarm!50}50.5 & \cellcolor{recwarm!56}56.3 & \cellcolor{recwarm!60}60.0 & \cellcolor{recwarm!59}59.1 & \cellcolor{recwarm!50}49.8 & \cellcolor{recwarm!52}52.5 \\
\bottomrule
\end{tabular}%
}

\vspace{2pt}
% {\footnotesize\raggedright
% Rung \textbf{TLV}, judged under the error-taxonomy rubric. The denominator is each column's
% scored questions (its $n$), which excludes OOM cells: those carry no answer, so they are
% neither a right answer nor a kind of wrong one. Categories are assigned first-match-wins in a
% fixed priority order, so they are mutually exclusive and every incorrect answer lands in
% exactly one; \emph{All incorrect} is their sum and equals 100 minus that column's accuracy.
% Bold is the largest value in each column within a group, the totals excepted.
% \textbf{The two groups are not comparable in level.} The answerable pool runs on oracle pages
% and the unanswerable pool on the bm25 $k{=}3$ page set, because an unanswerable question has
% no gold page to select. \emph{answered unanswerable} appears only in the second group: on a
% pool with an answer it can only be the judge misapplying the label, and on a pool without one
% it is the entire error, which is why \emph{All incorrect} reproduces it there.
% \emph{no answer reached} is a response that reasoned without committing, typically one that
% exhausted its decode budget mid-calculation; it is absent under the three modes that ask for
% a short answer. Read the answerable group beside Table~\ref{tab:prompting}, whose accuracies
% are what these rates decompose.\par}
\caption{\small Failure modes under six prompt instructions. Each cell is the share of that column's own question pool assigned the category, so a group's rows sum to its \emph{All incorrect} row rather than to 100.}
\label{tab:error_profile}
\end{table*}


% \subsection{Failure Modes by Evidence Source}
% \label{app:error_sources}
% Table~\ref{tab:error_source} is the tabular form of Figure~\ref{fig:error_sources}, with the same five small categories printed rather than pooled. Evidence-source columns overlap, since a question citing both a chart and a table is counted under each, so they do not partition the pool and \emph{All} pools every question once rather than summing across them. Each column is read against its own \emph{All incorrect} total, which is 100 minus that column's accuracy.
% % GENERATED by tables/scripts/ -- do not edit by hand, re-run the script.
% Numbers read from results/tables/analysis/*.csv.
% error_profile_source (results.errcat.jsonl)
% Shade ranges are fitted to the values present:
% Share of the source pool (\%): recwarm over [0.0, 57.6]
\definecolor{recwarm}{HTML}{D95926}
\begin{table*}[h]
\centering
\small
\setlength{\tabcolsep}{5pt}
\adjustbox{max width=\textwidth}{%
\begin{tabular}{l@{\hskip 14pt}rrrrrr}
\toprule
 & \multicolumn{6}{c}{\textbf{Share of the source pool (\%)}} \\
\cmidrule(lr){2-7}
\textbf{Error category} & \textbf{Text} & \textbf{Layout} & \textbf{Table} & \textbf{Chart} & \textbf{Figure} & \textbf{\emph{All}} \\
\midrule
insufficient evidence & \cellcolor{recwarm!8}7.3 & \cellcolor{recwarm!1}0.8 & \cellcolor{recwarm!6}6.2 & \cellcolor{recwarm!8}7.3 & \cellcolor{recwarm!8}7.6 & \cellcolor{recwarm!6}6.1 \\
no answer reached & \cellcolor{recwarm!4}3.8 & \cellcolor{recwarm!3}2.5 & \cellcolor{recwarm!15}14.4 & \cellcolor{recwarm!8}7.9 & \cellcolor{recwarm!5}4.5 & \cellcolor{recwarm!7}6.8 \\
equivocal & \cellcolor{recwarm!0}0.3 & \cellcolor{recwarm!1}0.8 & \cellcolor{recwarm!1}1.4 & \cellcolor{recwarm!1}1.1 & \cellcolor{recwarm!1}1.4 & \cellcolor{recwarm!1}1.1 \\
wrong count & \cellcolor{recwarm!13}\textbf{12.9} & \cellcolor{recwarm!30}\textbf{28.8} & \cellcolor{recwarm!8}7.7 & \cellcolor{recwarm!11}10.1 & \cellcolor{recwarm!25}\textbf{24.1} & \cellcolor{recwarm!15}\textbf{14.4} \\
wrong value & \cellcolor{recwarm!9}9.1 & \cellcolor{recwarm!6}5.9 & \cellcolor{recwarm!16}\textbf{15.4} & \cellcolor{recwarm!16}\textbf{15.7} & \cellcolor{recwarm!5}4.8 & \cellcolor{recwarm!9}9.0 \\
wrong entity & \cellcolor{recwarm!7}7.0 & \cellcolor{recwarm!7}6.8 & \cellcolor{recwarm!6}5.8 & \cellcolor{recwarm!11}10.7 & \cellcolor{recwarm!10}9.3 & \cellcolor{recwarm!8}7.7 \\
incomplete list & \cellcolor{recwarm!2}2.1 & \cellcolor{recwarm!2}1.7 & \cellcolor{recwarm!1}0.5 & \cellcolor{recwarm!2}1.7 & \cellcolor{recwarm!3}2.8 & \cellcolor{recwarm!2}1.8 \\
over-inclusive list & \cellcolor{recwarm!4}3.5 & \cellcolor{recwarm!3}2.5 & \cellcolor{recwarm!2}2.4 & \cellcolor{recwarm!0}0.0 & \cellcolor{recwarm!2}2.4 & \cellcolor{recwarm!2}2.2 \\
unit / format mismatch & \cellcolor{recwarm!3}2.8 & \cellcolor{recwarm!1}0.8 & \cellcolor{recwarm!1}1.4 & \cellcolor{recwarm!3}2.8 & \cellcolor{recwarm!0}0.3 & \cellcolor{recwarm!1}1.3 \\
other & \cellcolor{recwarm!0}0.0 & \cellcolor{recwarm!0}0.0 & \cellcolor{recwarm!0}0.0 & \cellcolor{recwarm!0}0.0 & \cellcolor{recwarm!0}0.3 & \cellcolor{recwarm!0}0.1 \\
\emph{All incorrect} & \cellcolor{recwarm!51}49.0 & \cellcolor{recwarm!53}50.8 & \cellcolor{recwarm!58}55.3 & \cellcolor{recwarm!60}57.3 & \cellcolor{recwarm!60}57.6 & \cellcolor{recwarm!53}50.5 \\
\bottomrule
\end{tabular}%
}

\vspace{2pt}
% {\footnotesize\raggedright
% Answerable pool, oracle pages, uninstructed prompt, rung \textbf{TLV}, judged under the
% error-taxonomy rubric. The denominator is each column's scored questions (its $n$), which
% excludes OOM cells. Categories are assigned first-match-wins in a fixed priority order, so
% every incorrect answer lands in exactly one and \emph{All incorrect} is their sum. Bold is
% the largest value in each column, the total excepted. \emph{answered unanswerable} is not
% shown: these questions have answers, so it can only be the judge misapplying the label.
% A question citing several sources is counted under each, so the source columns overlap and
% their $n$ do not sum to the pool; \emph{All} is the same questions pooled once each rather
% than a row-wise sum, which is why its $n$ is smaller than the columns beside it add to.
% \textbf{What changes across sources is which value is wrong, not whether the reasoner
% engages.} \emph{wrong count} takes 24.1\% of Figure questions and 28.8\% of Layout ones
% against 7.7\% of Table ones, and \emph{wrong value} runs the other way. The exception is
% \emph{no answer reached}, highest on Table at 14.4\%: a table is where the reasoner most
% often starts a calculation it does not finish. Table~\ref{tab:error_integration} splits the
% same pool by gold-page count, where that category is the whole story.\par}
\caption{\small Failure modes by the evidence a question cites. Each cell is the share of that source's own question pool assigned the category, so a column sums to its \emph{All incorrect} row rather than to 100.}
\label{tab:error_source}
\end{table*}


% \subsection{Failure Modes by Evidence-Page Count}
% \label{app:error_integration}
% Table~\ref{tab:error_integration} is the tabular form of Figure~\ref{fig:error_integration}, and prints the M$-$1 column that the figure draws as a strip: the difference between the pooled multi-page column and the single-page one, category by category. \emph{Multi} pools the two- and three-or-more-page columns rather than extending them, so the columns overlap and are never summed, and the 3+ column is thin enough to be read for trend rather than for level.
% % GENERATED by tables/scripts/ -- do not edit by hand, re-run the script.
% Numbers read from results/tables/analysis/*.csv.
% error_profile_hop (results.errcat.jsonl)
% Warm wash on the difference column deepens with its magnitude, in either direction.
\definecolor{gapred}{HTML}{B03A2E}

\begin{table}[h]
\centering
\small
\setlength{\tabcolsep}{5pt}
\renewcommand{\arraystretch}{1.12}
\begin{tabular}{@{}l rrr r@{\hskip 10pt} r@{}}
\toprule
 & \multicolumn{3}{c}{\textbf{Gold evidence pages}} & & \\
\cmidrule(lr){2-4}
\textbf{Error category} & \textbf{1} & \textbf{2} & \textbf{3+} & \textbf{Multi} & \textbf{M$-$1} \\
\midrule
insufficient evidence & 6.3 & 5.9 & 1.8 & 4.6 & \cellcolor{gapred!3}-1.7 \\
no answer reached & 2.5 & 15.1 & 8.1 & 12.9 & \cellcolor{gapred!16}+10.4 \\
equivocal & 0.6 & 2.1 & 0.9 & 1.7 & \cellcolor{gapred!2}+1.1 \\
wrong count & 10.1 & 10.9 & 39.6 & 20.0 & \cellcolor{gapred!15}+9.9 \\
wrong value & 7.4 & 14.6 & 4.5 & 11.4 & \cellcolor{gapred!6}+4.0 \\
wrong entity & 7.8 & 9.6 & 3.6 & 7.7 & \cellcolor{gapred!0}-0.1 \\
incomplete list & 1.7 & 1.7 & 2.7 & 2.0 & \cellcolor{gapred!0}+0.3 \\
over-inclusive list & 2.5 & 1.7 & 1.8 & 1.7 & \cellcolor{gapred!1}-0.8 \\
unit / format mismatch & 1.5 & 1.3 & 0.9 & 1.1 & \cellcolor{gapred!1}-0.4 \\
other & 0.2 & 0.0 & 0.0 & 0.0 & \cellcolor{gapred!0}-0.2 \\
\midrule
\emph{All incorrect} & 40.5 & 62.8 & 64.0 & 63.1 & \cellcolor{gapred!34}+22.6 \\
\bottomrule
\end{tabular}
\caption{Failure modes by the number of gold evidence pages. Each cell is the share of that column's own question pool assigned the category, so a column sums to its \emph{All incorrect} row rather than to 100.}
\label{tab:error_integration}
\end{table}

% Footnote, if the float needs one in the text:
% Answerable pool, oracle pages, uninstructed prompt, rung TLV, error-taxonomy rubric.
% `Multi` pools the 2-page and 3+-page columns, so the four columns overlap rather than
% partition and are not summed; M$-$1 is Multi minus the 1-page column, in points of
% failure rate. Questions with no annotated evidence page are dropped. The multi-page
% deficit is mostly a failure to finish: `no answer reached` runs 2.5% on one page
% against 12.9% across several, while the two categories meaning the reasoner read the
% evidence and misjudged it, `insufficient evidence` and `wrong entity`, barely move.

\clearpage
\section{Full Method Details}
\label{app:simple_method_details}